% --- CODIFICACIÓN E IDIOMA ---
\usepackage[utf8]{inputenc}
\usepackage[T1]{fontenc}
\usepackage[es-tabla, es-nodecimaldot]{babel}

% --- MATEMÁTICAS Y FÍSICA ---
\usepackage{mathtools}
\usepackage{amssymb}
\usepackage{amsthm}
\usepackage{physics}
\usepackage{bm}

% --- GRÁFICOS Y TABLAS ---
\usepackage{graphicx}
\graphicspath{{figures/}} % This allows you to just use {image.png} instead of {figures/image.png}
\usepackage{booktabs}
\usepackage{array}
\usepackage{tabularx}
\usepackage{ragged2e}
\renewcommand\tabularxcolumn[1]{m{#1}}

% --- CÓDIGO (Juia Support) ---
\usepackage{xcolor}
\usepackage{listings}
\lstdefinelanguage{Julia}%
  {morekeywords={abstract,break,case,catch,const,continue,do,else,elseif,%
      end,export,false,for,function,immutable,import,importall,if,in,%
      macro,module,otherwise,quote,return,switch,true,try,type,typealias,%
      using,while},%
   sensitive=true,%
   morecomment=[l]\#,%
   morestring=[b]",%
  }

% --- TEOREMAS (thmtools) ---
\usepackage{mdframed}
\usepackage{thmtools}


% 1. Estilo para Teoremas (Texto en cursiva, cabecera negrita sans-serif)
\declaretheoremstyle[
    headfont=\bfseries\sffamily, 
    bodyfont=\itshape,
    spaceabove=12pt, 
    spacebelow=12pt
]{miestilo-plano}

% 2. Estilo para Definiciones (Caja gris lateral, texto normal)
\declaretheoremstyle[
    headfont=\bfseries\sffamily,
    bodyfont=\normalfont,
    mdframed={
        linewidth=1pt,
        linecolor=gray!50,      % Color de la línea
        backgroundcolor=gray!5, % Fondo muy tenue
        topline=false, bottomline=false, rightline=false, % Solo barra izquierda
    }
]{miestilo-caja}

% --- DEFINICIÓN DE ENTORNOS ---
% Al ser 'article', numeramos por 'section'.
\declaretheorem[style=miestilo-plano, name=Teorema, numberwithin=chapter]{theorem}

% Estos comparten el contador con 'theorem' (e.g., Teorema 1.1, Lema 1.2)
\declaretheorem[style=miestilo-plano, name=Lema, sharenumber=theorem]{lemma}
\declaretheorem[style=miestilo-plano, name=Proposición, sharenumber=theorem]{proposition}
\declaretheorem[style=miestilo-plano, name=Corolario, sharenumber=theorem]{corollary}

% Definiciones y Ejemplos usan el estilo de caja
\declaretheorem[style=miestilo-caja, name=Definición, sharenumber=theorem]{definition}
\declaretheorem[style=miestilo-caja, name=Ejemplo, sharenumber=theorem]{example}

% Demostración
\makeatletter
\renewenvironment{proof}[1][\proofname]{\par
  \pushQED{\qed}
  \normalfont \topsep6\p@\@plus6\p@\relax
  \trivlist
  \item[\hskip\labelsep
        \bfseries\sffamily
        #1\@addpunct{.}]\ignorespaces
}{%
  \popQED\endtrivlist\@endpefalse
}
\makeatother


% --- HIPERVÍNCULOS ---
\usepackage[colorlinks=true, linkcolor=blue, citecolor=red, urlcolor=teal]{hyperref}
\usepackage[nameinlink, noabbrev]{cleveref}

% --- BIBLIOGRAFÍA ---
\usepackage[
    backend=biber, 
    style=apa,      % Opciones: apa, ieee, numeric, alphabetic
    sorting=nyt     % Ordenar por Nombre, Año, Título
]{biblatex}
\addbibresource{bib/referencias.bib} % Note the path change

% --- MÁRGENES ---
\usepackage{geometry}
\geometry{a4paper, total={170mm,257mm}, left=25mm, top=25mm}
